\chapter{A foundational multi-modal verifier}

\section{Framework}

Loom is a foundational verifier in that all its semantics and proofs are
expressed in a theorem prover and proof assistant (Lean). Therefore, in this
sense, final guarantees rely only on the Lean kernel. This shifts the burden of
trust from an external software or language like Dafny to a theorem prover,
thereby avoiding trusting large external verification pipelines.

To be precise, Loom is actually a framework for building multi-modal verifiers.
That is to say, it is not a single verifier but rather a meta-framework for
constructing many. But by abuse of language (and for convenience), we shall
simply refer to Loom as a verifier to mean one such verifier built upon Loom.
Later, we shall see Velvet which is an actual verifier built on top of Loom.

Loom provides some mathematical constructions and proofs of these constructions
expressed in Lean. Programs in Loom are modeled using monads. Monads are a way
to structure computations as a sequence of steps, where each step can also be
seen as a computation that potentially gives a value while also producing
potentially a side effect such as IO operation, failure, file reading or
writing. Monads give a compositional semantics. Programs become mathematical
objects inside Lean.

Loom uses an algebraic structure called monad transformer. Complex program
behaviors can be built using monad transformers. A monad transformer is like a
type constructor except it receives a monad and give back another monad. In our
case, we think of each monad transformer encoding some effect such as failure or
nondeterminism. By stacking monad transformers, one can then build programs with
multiple side effects. Verifiers can be constructed by stacking transformers.

\section{Shallow Embedding}

Lean is a very malleable language. Indeed, it is quite feasible to extend Lean
or build a DSL inside Lean. Generally speaking, there are two approaches to
building a language on the top of Lean. They are deep and shallow embedding.

In deep embedding, we use data structures (or inductive types in our case) to
build abstract syntax tree (AST) and write interpreters to interpret them. If we
go along this route to build our verifier language, we also have to verify the
correctness of the semantics of our interpreter.

Loom instead uses shallow embedding. Here programs are written directly in
Lean's language. No separate syntax or external language is needed. Benefits are
that it reuses Lean's type system and it now has direct access to proofs and
tactics. Programs are just Lean terms and this allows smooth integration with
theorem proving.


\section{Hybrid Workflow}

Another strong usage of Loom would be its hybrid workflow. It aims to seamlessly
blend automation via SMT solving with interactive theorem proving via Lean's
by-tactic blocks. This points to a unified framework for program execution,
verification and proof development in the same environment.

\section{Velvet}

We have mentioned before that Loom is not a single verifier but rather a
framework for building verifiers. We have been loosely referring to Loom as if
it were an actual verifier language but now we meet Velvet, an actual DSL and
verifier built on top of Lean.

The syntax of Velvet is inspired by Dafny. As with other verification languages,
Velvet is used to write programs with specifications. Programs are written in
imperative style (Velvet gives this imperative syntax). Velvet is also the name
of the monad in which all programs are written. Indeed, Velvet is an
abbreviation defined by

\begin{leanblock}
  abbrev VelvetM α := NonDetT DivM α
\end{leanblock}

That is, \lean{VelvetM} is a monad built from the divergence monad \lean{DivM}
(this encodes programs which can terminate or diverge, that is, to not
terminate) and the non-determinism monad transformer \lean{NonDetT} (this
encodes programs which can have non-deterministic behavior).

Indeed, all programs in Velvet DSL have the type \lean{VelvetM}.

Velvet combines automation (via \lean{loom_solve} which is an automation tactic
provided by Loom) and interactive proofs (which are just Lean tactics). More
automation is possible through backends SMT solvers such as Z3 or CVC5. This
gives a hybrid workflow in which verification usually proceeds by attacking with
automation and followed by, if needed, more interactive theorem proving
(tactics).

\section{Imperative Programming Syntax}

A majority of programmers are used to imperative-style programming like C or C++
or Java, and a lot of programming in current industry is mostly imperative.
Velvet gives an imperative-style syntax built on top of the functional language
Lean, made possible by leveraging the metaprogramming and elaboration techniques
of the language.

The following is a typical program written in Velvet.

\begin{leanblock}
method DoubleAll (mut arr: Array Int) return (u: Unit)
    ensures arr.size = arrOld.size
    ensures ∀ i, 0 ≤ i ∧ i < arr.size → arr[i]! = 2 * arrOld[i]!
    do
    let mut i := 0
    while i < arr.size
        invariant arr.size = arrOld.size
        invariant 0 ≤ i ∧ i ≤ arr.size
        invariant ∀ k, 0 ≤ k ∧ k < i → arr[k]! = 2 * arrOld[k]!
        invariant ∀ k, i ≤ k ∧ k < arr.size → arr[k]! = arrOld[k]!
        done_with i = arr.size
    do
        arr := arr.set! i (2 * arr[i]!)
        i := i + 1
    return
\end{leanblock}

The function starts with the keyword \lean{method}. The preconditions and
postconditions are provided using \lean{require} and \lean{ensures} clauses.
Multiple specifications are provided with multiple \lean{require} and
\lean{ensures} clauses. The \lean{while} loop requires invariants which are
provided using \lean{invariant} clause.

All functions written this way (using \lean{method}) are elaborated to standard
Lean programs or terms in the monad \lean{VelvetM}. Indeed, if we type
\lean{#print DoubleAll}, we see the output:

\begin{leanblock}
def DoubleAll : Array ℤ → VelvetM (Unit × Array ℤ)
\end{leanblock}

The above type can look a bit confusing (the extra \lean{Array ℤ} in the return)
but this is due to how mutable parameters are handled in Velvet which we will
see later.

\section{Multi-modal Verification}

To execute or run the function, we simply evaluate:

\begin{leanblock}
#eval (DoubleAll #[1, 2, 3]).run
\end{leanblock}

One can also perform property-based testing (PBT). This allows us to test
methods with random inputs before actual verification.

To prove the functional correctness and that the specifications hold for all
inputs, we do

\begin{leanblock}
prove_correct DoubleAll
  by loom_solve
\end{leanblock}

As mentioned above, \lean{loom_solve} is an automation tactic and usually forms
the first line of attack. For most simple proof obligations, Velvet claims, this
tactic is enough to finish the proof by itself. In the cases it does not, one
can try more tactics to finish the remaining goals.

By multi-modal, it refers to the fact that Velvet can handle program execution,
testing and deductive reasoning (formal verification) all in one environment.

\section{More on Loom and Velvet}

Velvet can do a lot more than what we have described in the above. For instance,
it allows programmers to change the semantics (that is from ``partial'' to
``total'' correctness for termination, or from ``demonic'' to ``angelic'' for
non-deterministic programs). These options are set in the beginning of the
program file. Usually, the options are set to ``partial correctness'' and
``demonic'' choice.

Switching from ``partial'' to ``total'' correctness restricts programmers to
write only terminating programs in the sense that when they write loops, for
instance, they also have to provide termination measures.

For more details, one can refer to this excellent documentation. [HERE]
